%%% Hlavní soubor. Zde se definují základní parametry a odkazuje se na ostatní části. %%%

%% Verze pro jednostranný tisk:
% Okraje: levý 40mm, pravý 25mm, horní a dolní 25mm
% (ale pozor, LaTeX si sám přidává 1in)
% \documentclass[12pt,a4paper]{report}
% \setlength\textwidth{145mm}
% \setlength\textheight{247mm}
% \setlength\oddsidemargin{15mm}
% \setlength\evensidemargin{15mm}
% \setlength\topmargin{0mm}
% \setlength\headsep{0mm}
% \setlength\headheight{0mm}
%% \openright zařídí, aby následující text začínal na pravé straně knihy
% \let\openright=\clearpage

% Pokud tiskneme oboustranně:
 \documentclass[12pt,a4paper,twoside,openright]{report}
 \setlength\textwidth{145mm}
 \setlength\textheight{247mm}
 \setlength\oddsidemargin{15mm}
 \setlength\evensidemargin{0mm}
 \setlength\topmargin{0mm}
 \setlength\headsep{0mm}
 \setlength\headheight{0mm}
 \let\openright=\cleardoublepage

%% Použité kódování znaků: obvykle latin2, cp1250 nebo utf8:
\usepackage[utf8]{inputenc}

%% Ostatní balíčky
\usepackage{graphicx}
\usepackage{amsthm}
\usepackage{color}


%% Balíček hyperref, kterým jdou vyrábět klikací odkazy v PDF,
%% ale hlavně ho používáme k uložení metadat do PDF (včetně obsahu).
%% POZOR, nezapomeňte vyplnit jméno práce a autora.
%\usepackage[ps2pdf,unicode]{hyperref}   % Musí být za všemi ostatními balíčky
%\hypersetup{pdftitle=: Algorithms for solving strong positional games with high symmetry}
%\hypersetup{pdfauthor= Jitka Novotná}

%%% Drobné úpravy stylu

% Tato makra přesvědčují mírně ošklivým trikem LaTeX, aby hlavičky kapitol
% sázel příčetněji a nevynechával nad nimi spoustu místa. Směle ignorujte.
\makeatletter
\def\@makechapterhead#1{
  {\parindent \z@ \raggedright \normalfont
   \Huge\bfseries \thechapter. #1
   \par\nobreak
   \vskip 20\p@
}}
\def\@makeschapterhead#1{
  {\parindent \z@ \raggedright \normalfont
   \Huge\bfseries #1
   \par\nobreak
   \vskip 20\p@
}}
\makeatother

% Toto makro definuje kapitolu, která není očíslovaná, ale je uvedena v obsahu.
\def\chapwithtoc#1{
\chapter*{#1}
\addcontentsline{toc}{chapter}{#1}
}

% Toto makro slouží k poznámkám, text se zobrazi červeně TODO
\def\TODO#1{
\textcolor{red}{#1}
}
\def\obr#1{
\textcolor{red}{#1}
}
\let \node \textsf
\let \value \textsf
\let \fce \textsl
\let \sec \textsc

\begin{document}

% Trochu volnější nastavení dělení slov, než je default.
\lefthyphenmin=2
\righthyphenmin=2

%%% Titulní strana práce

\pagestyle{empty}
\begin{center}

\large

Charles University in Prague

\medskip

Faculty of Mathematics and Physics

\vfill

{\bf\Large BACHELOR THESIS}

\vfill

\centerline{\mbox{\includegraphics[width=60mm]{../img/logo.eps}}}

\vfill
\vspace{5mm}

{\LARGE Jitka Novotná}

\vspace{15mm}

% Název práce přesně podle zadání
{\LARGE\bfseries Algorithms for solving strong positional games with high symmetry}

\vfill

% Název katedry nebo ústavu, kde byla práce oficiálně zadána
% (dle Organizační struktury MFF UK)
Computer Science Institute of Charles University
\vfill

\begin{tabular}{rl}

Supervisor of the bachelor thesis: & Tomáš Valla \\
\noalign{\vspace{2mm}}
Study programme: & Informatika \\
\noalign{\vspace{2mm}}
Specialization: & Obecná informatika \\
\end{tabular}

\vfill

% Zde doplňte rok
Prague 2012

\end{center}

\newpage

%%% Následuje vevázaný list -- kopie podepsaného "Zadání bakalářské práce".
%%% Toto zadání NENÍ součástí elektronické verze práce, nescanovat.

%%% Na tomto místě mohou být napsána případná poděkování (vedoucímu práce,
%%% konzultantovi, tomu, kdo zapůjčil software, literaturu apod.)

\openright

\noindent
Dedication.
\TODO{ Tom, rodiče, Roman - laděni, Tom - vše, Honza - podpora +?angličtina, 
Paulie -předvyber článku, Pavel to s Cachy, Brum - chytre rady}

\newpage

%%% Strana s čestným prohlášením k bakalářské práci

\vglue 0pt plus 1fill

\noindent
I declare that I carried out this bachelor thesis independently, and only with the cited
sources, literature and other professional sources.

\medskip\noindent
I understand that my work relates to the rights and obligations under the Act No.
121/2000 Coll., the Copyright Act, as amended, in particular the fact that the Charles
University in Prague has the right to conclude a license agreement on the use of this
work as a school work pursuant to Section 60 paragraph 1 of the Copyright Act.

\vspace{10mm}

\hbox{\hbox to 0.5\hsize{%
In ........ date ............
\hss}\hbox to 0.5\hsize{%
signature of the author \TODO{ Tady prijde jen rucni nebo i vytiskly podpis? }
\hss}}

\vspace{20mm}
\newpage

%%% Povinná informační strana bakalářské práce

\vbox to 0.5\vsize{
\setlength\parindent{0mm}
\setlength\parskip{5mm}

Název práce:
Algorithms for solving strong positional games with high symmetry
% přesně dle zadání

Autor:
Jitka Novotná

Katedra:  % Případně Ústav:
% Název katedry či ústavu, kde byla práce oficiálně zadána
Informatický ústav Univerzity Karlovy
% dle Organizační struktury MFF UK

Vedoucí bakalářské práce:
RNDr. Tomáš Valla, Informatický ústav Univerzity Karlovy \TODO{je to pracoviste dobře?}
% Jméno a příjmení s tituly, pracoviště
% dle Organizační struktury MFF UK, případně plný název pracoviště mimo MFF UK

Abstrakt:
% abstrakt v rozsahu 80-200 slov; nejedná se však o opis zadání bakalářské práce
\TODO{TODO + zkontrolovat jestli sedi s anglictinou}

Klíčová slova:
% 3 až 5 klíčových slov
\TODO{TODO PN-search, silné poziční hry, kliková hra, vysoká hra, vylepšení PN-searche}

\vss}\nobreak\vbox to 0.49\vsize{
\setlength\parindent{0mm}
\setlength\parskip{5mm}

Title:
Algorithms for solving strong positional games with high symmetry
% přesný překlad názvu práce v angličtině

Author:
Jitka Novotná

Department:
Computer Science Institute of Charles University
% Název katedry či ústavu, kde byla práce oficiálně zadána
% dle Organizační struktury MFF UK v angličtině

Supervisor:
RNDr. Tomáš Valla, Informatický ústav Univerzity Karlovy 
% dle Organizační struktury MFF UK, případně plný název pracoviště
% mimo MFF UK v angličtině

Abstract:
% abstrakt v rozsahu 80-200 slov v angličtině; nejedná se však o překlad
% zadání bakalářské práce
\TODO{TODO}

Keywords:
% 3 až 5 klíčových slov v angličtině

\vss}

\newpage

%%% Strana s automaticky generovaným obsahem bakalářské práce. U matematických
%%% prací je přípustné, aby seznam tabulek a zkratek, existují-li, byl umístěn
%%% na začátku práce, místo na jejím konci.

%\openright
\pagestyle{plain}
\setcounter{page}{1}
\tableofcontents

%%% Jednotlivé kapitoly práce jsou pro přehlednost uloženy v samostatných souborech
\chapter*{Introduction}
\addcontentsline{toc}{chapter}{Introduction}

This thesis introduces an algorithm for solving strong positional games
with high symmetry. We describe the most important algorithm for solving 
combinatorial games -- PN-search and its various enhancements which change
it so much that they can be considered new algorithms. 

%First, we present several known enhancements. Second, some enhancements
%which we designed. Then we test them on the example of the clique game.
%
%Our solver proved that in (6,4) and (7,4)-clique game the second player has a draw
%strategy. This was expected but not proved yet.

\medskip

PN-search was introduced by Allis~\cite{Allis} in 1994. It has become well know algorithm
and basis for a lot of further work. It was successfully applied for a wide class
of games and many enhancements was invited. We introduce some of them and how they can be useful
in case of a game with high symmetry. We consider which  enhancements
can be used together and what problems can arise in combinations.
We also introduce some enhancements which we discovered independently.

\medskip

We do not determine neither asymptotic time complexity nor asymptotic space complexity of the algorithms in this thesis.
It is not clear what time/space complexity means in games solving -- 
if the input of the  algorithm are rules of the game most of games can be solved in constant time.

On the other hand when the input of the algorithm is a game state
of a positional game with $n$ positions solving the game is at least hard problem.
($n = {N \choose 2}$ in ($N$,$K$)-clique game). The number of game states in
such a game is \TODO{bigO} $O(2^{N \choose 2})$. It is not necessary to visit them all,
however, the number of the game states we must visit even for proving that a strategy
is winning is also $O(2^{N \choose 2})$. From this view all improvements in this
thesis are just heuristic.

The size of the partial game tree which we must create for solving the game is 
much more important then asymptotic complexity.

\subsubsection{Outline}

In the first chapter we informally introduce some combinatorial
games terminology. We suppose the reader has general knowledge about graph theory, 
generally see Deisel \ref{TODO}. In the first chapter we also define clique game and give
solutions of some cases of a clique game.

Known algorithms and their enhancements are introduced in the second chapter. Usually
algorithms are designed for general combinatorial game. We apply them to
strong positional games, which simplifies the descriptions. We start with a meta
algorithm of game solving. Some specific algorithms are described after. We choose PN-search to
implement and to be described in detail.

Theoretical results which we found independently are described in the third chapter.
In subsections \ref{} to \ref{} TODO are enhancements of PN-search. Then we introduce some rules for
deleting from cache which are designed for strong positional games. We consider
a problem which can occur when we combine some enhancement and we
propose solution for the problem. In the end of the chapter we focus on
a graph of a clique game \TODO{toto do dreti kap}. We describe several ways of storing
the game state and also how to normalize them.

The last chapter introduces our software and presents experimental results.
First, we describe the partial game tree created by our solver during solving
(6,4) and (7,4)-clique games. Then we compare some mentioned enhancements. We
study if they help to solve clique game. In the end of the chapter we describe in
more detail the software implementation.

\chapter{Games}

\TODO{node, 
game position
game tree=DAG 
free winning set,
symmetric game 
finite
hrozba
definice turn-to co bezne, to co ja,
free edge/claim 
grafem myslim dvoubarevny neorientovany
ze druhy nema vyhravajici strategii
definice strategie
win,draw position
free node
} 
\section{ Definition }

\subsection{ Two players sequential game}

\subsection{ Positional game}
\TODO{ala piskvorky}

\subsubsection{ Strong game }

\subsubsection{ Maker-Breaker games }

\subsection{ Other property }
\TODO{ Hight symmetry, strong branching factor}

\section{ Clique game }

\subsection{ Rules }

\subsection{ Ramsey theorem }

\subsection{ Solution for some N, K }

\subsection{ Space and time complexity }
\TODO{
neresim slozitost ale velikost stromu 
udelat bubu jak je ten problem tezky
rict ze nema smysl resit asimptotickou slozitost fci v N ze na konstantach a heuristikach zalezi vic
}

 

\section{ Other games }
\TODO{treba ty ctverecky, mb varianta klikvo hry vysledky, citace, reversni varianta}


\chapter{Algorithms and enhancements}

\section{Overview of algorithms}

\subsection{Meta algorithm}

\TODO{hrana (u,v)}

We will describe a~meta algorithm how to prove that a finite two-player
positional game is determined, that is there is a~winning strategy for the first
player or draw strategy for the second.

All game states can be represented as a~rooted tree. Each game position is a
node and if someone can turn from one position to another, there is an edge
between the corresponding nodes. We will not try to build a~complete game tree,
representing every possible position in the game, because it would be too large
to be represented in memory. Instead, we will incrementally build a partial game 
tree starting from the empty tree, trying to keep the partial game tree as small as
possible.

We assign to each node one of three possible values -- \value{true},
\value{false}, \value{unknown}. Node will have value \value{true} when we prove
that there is a~winning strategy for the first player and value \value{false}
when we prove that there is a~draw strategy for the second player. Otherwise
the node has value \value{unknown}. We say that node is solved if it has value
\value{true} or \value{false}.

The algorithm starts with one node with value \value{unknown}. In each iteration the
algorithm does one of the two things:

\begin{itemize} 
	\item{Develop node} We choose one leaf node with value \value{unknown}
		and develop it, which means that we create its sons. When 
		a son is created we evaluate it (we check whether one of the players
		won).
	\item{Update node} We choose one developed node with value \value{unknown}
		and look over its sons. We check if we can assign value \value{true} 
		or \value{false} to the node.
\end{itemize}

Algorithm ends when root is solved and so game is determined.

There are two types of nodes. When the first player is on turn it is an \node{or}
node. To assign value \value{true} for it, it is enough if one of its sons
have value \value{true}. Then, the first player has a winning strategy which starts by the
turn to the node with value \value{true}. On the other hand, when the first player has
no winning strategy after any turn he has no winning strategy at all. So we can
assign value \value{false} when all sons have value \value{false}. If none of above holds
the value of the node remains \value{unknown}. Hence value of an \node{or} node can be found by
an operation which is very similar to simple logical \fce{or}.

The second type of node is an \node{and} node. The second player is on turn. Assigning values
works in a similar way, we just use \fce{and}-like operation instead of \fce{or}-like.

\obr{obrazek 1) and kolecko} 

\subsection{Depth-first search} We can use ordinary depth-first search for
searching the complete game tree and solving the root.  We start by creating
the root. We develop root and choose one of it sons. We solve the son
recursively and update the root. We continue solving the sons and updating
until the root is solved.

This algorithm is a suitable solution for the clique game. The depth of the
tree is $N \choose 2$ so it doesn't take much memory. This is obviously slow
but if we use some of enhancements of pn-search which are described later, this
algorithm can find a solution for $N=5$ and maybe more.

\subsection{Alpha-beta}

Alpha-beta is a well known algorithm for finding good strategies. \TODO{citace}.
We can use it for any node $n$. It searches a part of the subtree under $n$ and for
nodes in some depth uses a rating function and returns the rating of $n$. The rating
which alpha-beta returns can be used for deciding to which position to turn. It
is useful as a~heuristic and can be used for example in AI and also in our
problem. We can use depth-first search and when we make a decision which sons
to solve, we first rate all its sons by alpha-beta and then solve them in order
determined by the ratings.

\subsection{PN-search}

PN-search was introduced by Allis \TODO{citace}. We will describe an immediate evaluation
variant of it which is better suited for our problem, because a test if a position is winning is
is quick.

PN-search is a~best-first search. We need to specify which node is the best for our problem.
 We start with some definitions, which were used in \TODO{allis}. 

\newtheorem*{prove}{Definition}	
\begin{prove}
We will say that we \emph{prove} a node if we proved that the first player has winning strategy
from it and we can assign \value{true} to it. 
\end{prove}

\newtheorem*{disprove}{Definition}	
\begin{disprove}
We will say that we \emph{disprove} node if we proved that the second player has draw strategy
from it and we can assign \value{false} to it. 
\end{disprove}

\newtheorem*{proofSet}{Definition}	
\begin{proofSet}
	For any partial game tree $T$ a~set of leaf nodes with value \value{unknown} $S$ is a~{\sl proof set}
	if proving all nodes within $S$ proves $T$.
\end{proofSet}

\newtheorem*{disproofSet}{Definition}	
\begin{disproofSet}
	For any partial game tree $T$ a~set of leaf nodes with value \value{unknown} $S$ is a~{\sl disproof set}
	if disproving all nodes within $S$ disproves $T$.
\end{disproofSet}

\newtheorem*{proofNumber}{Definition} 
\begin{proofNumber}
	For any partial game tree $T$, the {\sl proof number} of $T$ is defined as the 
	cardinality of the smallest proof set of $T$.
\end{proofNumber}

\newtheorem*{disproofNumber}{Definition}	
\begin{disproofNumber}
	For any partial game tree $T$, the {\sl disproof number} of $T$ is defined as the 
	cardinality of the smallest disproof set of $T$.
\end{disproofNumber}

We will show how proof and disproof numbers work and how they can be calculated
on some examples. At picture \obr{2)} you can see a tree with proof numbers.
Nodes with value \value{true} have proof number 0 --- there is nothing left to
prove. Nodes with value \value{false} have proof number $ \infty $ --- there
exists no set of nodes which proves it. Leaf nodes with value \value{unknown}
have proof number 1 because it is enough to prove the node itself. Internal
\node{and} node has proof number which is equal to sum of the proof numbers of
its sons, because for proving the \node{and} node we must prove all its sons by
proving their proof sets. Internal \node{or} node has proof number equal to
minimum of proof numbers of it sons because we need to prove at least one of
its son.

Determining the disproof number works in a similar way. As you can see an picture
\obr{3)}. 

\newtheorem*{mostProvingNode}{Definition}	
\begin{mostProvingNode}
	For any partial game tree $T$ a~{\sl most proving node} of $T$ is a~node, which by 
	obtaining the value \value{true} reduces $T$'s proof number by 1, while by obtaining the
	value \value{false} reduces $T$'s disproof number by 1.
\end{mostProvingNode}

Most-proving node is a node which helps solve the root in both cases (proving or disproving it)
so we will use 
it as the best node in best-first search. We will show that such a node exists in every
partial tree.

\begin{lemma}
	In every partial game tree there is a most proving node.
\end{lemma}

\begin{proof}

	We will prove it by induction on depth of partial game tree.
	It is enough to show that there is a minimal proof set and a minimal disproof set having a 
	common node. Proving the node reduces the proof set and disproving it reduces the
	disproof set.

	\begin{itemize} 
		\item{Basic} 
			If root is a~leaf node, both the proving and the disproving set 
			contain just the root.
		\item{Induction step}
			Let the root be an \node{and} node. Minimal proof set of the root is the minimal proof set
			of its son $s$ with minimal proof number. Minimal disproof set of the root
			is the union of minimal disproof sets of its sons, so it contains
			the minimal disproof set of $s$.

			Imagine a~game (sub)tree where $s$ is root. By induction, his
			minimal proof set and minimal disproof set have a common node, which is also 
			in the intersection of minimal proof and minimal disproof sets of the original root and
			is the most proving node.

			Proof for \node{or} root proceeds analogously.
	\end{itemize}
\end{proof}

Now we informally describe the whole pn-search, as detailed in pseudocode in
\TODO{pseudkod}. 

We store a partial game tree and proof and disproof numbers of each node. We
start with one node --- the root. Then we repeat the following three steps
until the root is solved. 

The first step is \fce{selectMostProving(node)}. This function works similarly
to the proof that a most proving node exists. \obr{4)obrazek} We start in the
root and continue by choosing the son which has the minimal proof number when
we are in \node{and} node, respectively the minimal  disprove number in
\node{or} node, until we are in a~leaf node. When there are more sons with a
minimal number we choose the leftmost one (we fix an arbitrary ordering). 

The second step is \fce{developNode(mostProvingNode)}. There we create sons of
the \fce{mostProvingNode} and evaluate them.

The third step is \fce{updateAncestors(mostProvingNode)}. Proof and disproof
numbers in partial game tree have changed so we must update their stored
values. We start updating \fce{mostProvingNode} and continue by updating its
ancestors up to the root. At the end of these three steps the (dis)proof numbers
are again correct, because only numbers above \fce{mostProvingNode} could have
changed.

\TODO{opsat pseudokod z allis}

\subsection{Choice of algorithm}

\TODO{prepsat}
ze tetusim jestli je vyhr/prohr u nejakych nahodnych nodu
We choose pn-search to our problem. It seems to be quicker then alpha-beta. \TODO{citace} Alpha-beta
could be part of pn-search when we use rating from it as a~heuristic (\TODO{11}, \ref{ord}).
And finally there are many enhancements and heuristics which seem to help reduce the partial game tree size  
then the choice the algorithm.

\section{Enhancements of PN-search}

\TODO{obecny kec ze to
	dela prakticky jine algoritmy 
	a ze muze byt problem s kombinovanim,
	ze vetsina byla v allise at nemusim citovat furt,
ze pd-numbers MPN nejsou podle definice
ze to aplikuju hlavne na ten zakladny
ze tady hlavne predstavuju jak funguji neresim jak moc pomuzou
}

\subsection{Path from root}

Proof number search is best-first search. Its disadvantage is that searching for most-proving node takes
a lot of time. 

\newtheorem*{currentPath}{Definition}	
\begin{currentPath}
	A~{\sl current path} is a~path from root to the most proving node which is
	used when we search most proving node. \TODO{ktera cesta}
\end{currentPath}

In pn-search each iteration starts at the root and descends down. After developing
most proving node we update proof and disproof numbers starting in most-proving and 
return the way back into root. So we traverse the current path twice. 

There are two enhancements which reduce the number of nodes traversed to select the
most proving node. 

\subsubsection{Last changed node} \label{last}

We make two basic observations.
\begin{itemize}
\item When there is no change during node update, traversal can be stop even though we
aren't at root.
\item Current paths of two consecutive iterations have the same beginning. They can differ
only in part where proof or disproof number was updated.
\end{itemize}

\newtheorem*{currentNode}{Definition}	
\begin{currentNode}
A~{\sl current node} is a~node from current path. It is the lowest unchanged node or
root if everything was changed.
\end{currentNode}

Enhanced pn-search works like ordinary pn-search. It has one variable
called $currentNode$, at the beginning value of $currentNode$ is root. Each
$selectMostProvingNode()$ starts in $currentNode$. And each $updateAncestors()$
ends at the node where proof and disproof numbers haven't been changed and sets 
$currentNode$ to the last changed node. You can see as it works there. \obr{(5))}

\subsubsection{DF-PN search} \label{dfpn}

We can save even more traversing. In two consecutive iterations current paths
have common beginning. \obr{6.1)} We reduced the part when proof and disproof
numbers stay same. Now it's time to reduce part when they are changing and work
as in \obr{6.2b)}. We need to know the highest node which will not be in new
current path. We use thresholds for it --- we will store two new numbers for
each node on the current path. They will be set so that when (dis)proof number
will be smaller than its threshold we stay in the same subtree. If its equal or
bigger we must either update also its parent or the son with minimal (dis)proof
number will by different. 

Let $p$ and $d$ be proof and disproof number of node $n$. Let $pt$ and $dt$ be
their thresholds. When $n$ is \node{and} node we assume that its sons $1\ldots k$
are ordered so $p_1 < p_2 < \ldots < p_k$. Rules for setting new thresholds
were deduced in \TODO{cit: epsilon trick str4}:\obr{6.3)} 
\begin{eqnarray*} 
	pt_1 = min(pt, p_2+1), dt_1 = dt - d + d_1.
\end{eqnarray*}

When $n$ is \node{or} node we assume that $d_1 < d_2 < \ldots < d_k$. Rules
are symmetrical:  
\begin{eqnarray*} 
	pt_1 = pt-p+p_1, dt_1 = min(dt,d_2+1).
\end{eqnarray*}

When we delay updateAncestors to the time when $p > pt$ or $d > dt$ we can make
another changes. 

Main point is that we don't need to store whole partial game three. It is
enough to store nodes from current path. Algorithm could be implemented
recursively (but we didn't done it) and for this reason it is called depth-first.
It isn't typical depth-first search because one node can be visited more
then once \TODO{je tu potreba priklid nebo to je videt?}. 

DP-PN search is typically implemented with cache used for storing as many
nodes from partial game tree as possible. 

\subsection{Transpositions in DAG} \label{DAG}
 
As we see in picture \obr{8) 3 obr,neco jen ... misto hran, nespojene, spojene
jen poradi, spojene i symetrie popisek..} of chess tree we are exploring one
game state many times. Basic enhancement is to join that nodes into one.

The problem is that we have game DAG \footnote{There aren't cycles because in a
positional game all edges leads into position with more positions occupied.}
instead of game tree and proof and disproof numbers cannot by computed like
above because one node can increase proof number of another node more than by
one \obr{7)obr Proof and disproof numbers a) by definition --- disproof set for
root are three undeveloped nodes b) computed as
above}. However $proofNumbers$ and $disproofNumbers$ computed by
PN-search can be useful even though they can be higher then proof and disproof
number as they was defined. 

We can use PN-search with a~few modifications. When we generate children we
check whether they already exist. We can use hash table for this. Second
modification is that we need to update all parents. This could cause problems
if used together with other enhancements as last changed node \ref{last} but if
we use ordinary PN-search it works. Allis proof it in \cite{allis} (page
39--40).

One game state in turn $t$ can be visited in the worst case $t!$ times because
positions can be occupied in any order. So joining them is a~good idea and
there exist many enhancements of this enhancement.

In this thesis we will try to use advantage of symmetry as much as possible.
So we want in addition join game positions which are isomorph. See \obr{9) }
how it works on tic-tac-toe game tree. However in clique game we cannot join
all isomorphuos game position into one because finding \TODO{jak
se to pise?} represent is hard problem. So we will use ``almost
representative graph'' \TODO{novy nazev} instead. See section \ref{norm}
for definition.

\textbf{Note}: We will use ``game tree'' even if we talk about game DAG. It can be little confusing 
but in many cases there isn't any reason to distinguish if we talk about tree or DAG. 
We will use ``leaf node'' when we are talking about undeveloped node. 

%\subsection{Heuristic 1 1}
%%
%\TOD{napsat co je me?}
%
%By definition the proof and disproof numbers of a~leaf node with value
%\value{unknown} are 1 and 1. In general and/or tree we don't now anything more
%and we must use this assignment, when we know more we can set these numbers
%differently and break definition.
%
%The way the algorithm proceeds can seems similar to breath-first search when we run it
%on some tree where up to some depth initially all nodes have value \value{unknown}
%and each node on same level has same number of sons. \ODO{obrazek vyvoj}
%
%For this reason it is useful to initially set proof and disproof numbers
%differently. The first we can use the numbers which the node will probably have if
%we develop it. For example if the \node{and} node has $n$ sons we give it proof
%number $n$ and disproof number $1$.
%
%We will call this enhancement $n 1$. It causes the algorithm will develop left most
%nodes until it finds some node with value \value{true} or \value{false}. \ODO{obrazek}
%
%\ODO{!!!!!!!!!!!!!!!!!!!!je tohle dobre? ? myslim ze je}
%There is other way how to assign the numbers. \ODO{citace} When PN-search works on partial
%tree where are only few or no value \value{true} and \value{false} nodes, it prefers node with 
%fewer sons. We can initially set $n$ using some heuristic. When it is \node{and} node with
%$n$ sons and position is good for the first player we give it numbers $m 1$ ($m$ is slightly less
%then $n$). When it is \node{or} node with $n$ sons and position is bad for the first player we
%give it numbers $1 m$ ($m$ is slightly more then $n$).

\subsection{Weak PN-search} \label{weak}

When we use transposition into DAG one node can by counted in (dis)proof number many times.
Hence we show an enhancement which tries to minimize this disadvantage. 

\TODO{tady mam v kodu nejakou vyjimku to je pro 0 nebo pro 1?}
\TODO{zkontrolovat ze je tam ta \^2 jen jednou}
\TODO{da se to cele nejak jednoduse vysvetlit? \^2.. kdyz vyvracim zalezi na branching factor kdyz dokazuju prehodit}
We will count proof and disproof number in a~different way:
\begin{enumerate} 
	\item Let $n$ be a~leaf node. It depends on value of node.
		\begin{itemize}
			\item \value{true} $p(n)=0$, $d(n)=\infty$
			\item \value{false} $p(n)=\infty$, $d(n)=0$
			\item \value{unknown} $p(n)=1$, $d(n)=1$
		\end{itemize}
	\item Let $n$ be an internal \node{or} node with $k$ sons. Then \newline
		$p(n) = min_{1 \le i \le k}(n_i)$, \newline
		$d(n) = max_{1 \le i \le k}(n_i) + (k-1)^2$. 
	\item Let $n$ be an internal \node{and} node with $k$ sons. Then \newline
		$p(n) = max_{1 \le i \le k}(n_i) + (k-1)$, \newline 
		$d(n) = min_{1 \le i \le k}(n_i)$.
\end{enumerate}
		
\subsection{Deleting solved subtrees} \label{delete}

To save memory we can delete node after it is solved. However this doesn't 
help --- see \TODO{4.kapitola}

\subsection{ Prohrano pokud neni vyhravajici linie}
\TODO{dopsat az po prvni kapitole}

\subsection{PN-set}
\TODO{dopsat az napragramuju}

\subsection{1+\TODO{epsilon} trick}
\TODO{dopsat az napragramuju}


\TODO{tady mozna PN2, DB}


\chapter{Further theoretical improvements}

In this chapter we examine various improvements to pn-search and clique
game. All results
were independently discovered by the author of this thesis.

\section{Other enhancements of PN-search} 
Here we describe several enhancements
and also show some solutions for problems which arise when we combine several
\sec{PN-search} enhancements together. 

\subsection{Ordering of sons} \label{ord}

The ordering of sons of some node could be important. When there are two or more
sons with a~minimal (dis)proof number we usually choose the leftmost but there could
be better variant so we order sons by some heuristic.

\subsection{Deleting solved sons}

When we disproof a~son of \node{or} node (and symmetrically when we proof a~son
of \node{and} node) we do not solve the node but we can we can delete the edge
between them. The node will have less sons so the sum or max or min of numbers
could be counted faster and update ancestors will be also faster. A~bigger
improvement is when using the \sec{weak PN-search}. \sec{Weak PN-search} depends on
the branching factor, which decreases.

\subsection{Deleting parallel edges } \label{multi-edges} 

When we use \sec{transposition to DAG} multiple parallel edges may be
created, for example on the beginning of clique game. The root has $N \choose
2$ sons. Sons, all of them represent $K_N$ with one edge colored. When we
normalise them they will all be the same, and so we have $N \choose 2$ parallel
edge from the root to the same node. For a~reason similar to the previous
subsection, it is a~good idea to change parallel edges to a~single edge.

\subsection{Developing a~node without creating sons} \label{lastenh}

This enhancement saves memory. The partial game tree is largest when the root is
solved. In this state there is a~lot of nodes which are undeveloped. \obr{ 11) vysvetlivky
hotovy strom se spoustou listu} We could reduce memory requirements if we do not
create part of them. 

When we are developing a~node $n$ we create and evaluate all $n$ sons in order to
get proof and disproof numbers, as usually. Afterwards, we delete these sons. When
we find that the  most proving node is a~son of $n$ we create all sons of $n$ again and
now let them stay in the partial game tree. Every node is created at most twice so
time requirements rise at most twice. \obr{12) 11) po tomhle} 

Perhaps we can save even more memory when we use some counter in nodes
and save only the developed sons or all the sons if the percentage or number 
of developed sons is above some threshold.

\subsection{Cache}

There is a~possibility to use a node cache, instead of the representation of the whole
tree, when we use \sec{DF-PN-search}.
We can use a cache also when we use ordinary \sec{PN-search} but we must
do some modifications in this case: 

\begin{enumerate}
\item Do not delete nodes on the current path.
\item Stop \fce{updateAncestors()} when an ancestor is missing from the cache.
\item When a son is missing in the cache, create it again.
\end{enumerate}

Both \sec{DF-PN-search} and modified \sec{PN-search} perform poorly when the
cache is really small (slightly more then the maximal current path length). In
this case it could happen that algorithm never solves the root. It could happen
that the root has two sons $s_1$ and $s_2$  and algorithm runs for a~while in
the subtree of $s_1$ and then in the subtree of $s_2$, deleting all the nodes
of the subtree of $s_1$ from the cache. When it returns back below $s_1$ to
improve it, it cannot because it must create the nodes again (hence deleting
the nodes from subtree of $s_2$) and the situation repeats ad infinitum.
Hopefully this does not happen with a big cache.

Here we specify how to decide which node to delete from the cache. We cannot
delete a node from the current path and sons of a node which is just being
developed. We can set other rules about which nodes cannot be deleted or set
priorities for deleting. This can have significant impact on the performance of
the algorithm as we might delete results which took a long time to compute.

When we set too many rules it can happen that there is no node which can be
deleted and in this case we halt whole program. In our experience with our
solver is that it happens very soon after program starts or never.

\subsubsection{List of possible rules for deleting from cache:}
\begin{itemize}
\item Do not delete a node on the current path and sons of a node which is being developed.
\item Do not prefer nodes with value \val{true} or \val{false} -- they are 
      more important than the nodes with value \val{unknown} (but we need a really big
      cache for it).
\item Do not prefer nodes with many parents -- these are more important than nodes
      with few parents.
\item Prefer nodes which are not developed -- we do not lose any important information
      when we delete these. 
\end{itemize}

\subsection{Variants of \fce{updateAncestors()}}

There are three enhancements which influence the function \fce{updateAncestors()}. 

The first is using \sec{last changed node} to shorten the path. 

The second is \sec{DF-PN-search} (even if there is no function named \fce{updateAncestors()} and it is
updating lazily).

The third is \sec{transposition to DAG}. We found no article
considering changes in \fce{updateAncestors()} when node has more then one parent.

When we use the ordinary \fce{updateAncestors()} together with
\sec{transposition to DAG} and one of the other two enhancements,
\TODO{pravopis} it could happen that some numbers are changed on the current path
above an unchanged node. For an example, see Figure \obr{10)} root
was changed and its son on current path no. It could also happen that the most
proving node found by the algorithm could be different than most proving node
by the definition, see Figure \obr{15)}. This is not as big problem as it seems to be -- the
algorithm is still correct. An unusual situation can happen when we also use
\sec{deleting solved subtrees} -- as a~Czech proverb warns we could ``saw off
a~branch which is holding us''
 
To avoid it we must use new \fce{updateAncestors()} which updates \emph{all}
paths up to root, and better \fce{selectMostProving()} as we show in the next
subsection. 

\subsubsection{Simple updateAncestors()}

A~simple way how to define such \fce{updateAncestors()} is recursively. The
function updates a node and if it changed either proof or disproof
number, we call the function recursively on parents. It could happen that
some nodes are updated many times. So its parents are also updated many times
and the number of updates is growing with the number of possible paths.

\subsubsection{Level updateAncestors()}

A~better solution is to update each node on some level and create a~list of
parents which need to be updated. After we dispose of duplicates in the list,
we do the same on a~higher level. 

\subsubsection{One-visit updateAncestors()}

Another solution is to store more information, namely in which iteration of
\sec{PN-search} was the node last updated. In this case we can update
recursively and stop when we would update some node twice. It could happen
that this algorithm does not count right values but we can still use it as an effective 
approximation.

\subsection{Variants of \fce{selectMostProving()}}

As you can see in Figure \obr{10}, it could happen that some numbers are changed on
the current path above an unchanged node. So the observation from Section \sec{last changed
node} (\ref{last}) does not hold and it can happen that the algorithm select a most
proving node wrongly. There is an~example in Figure \obr{15}. We found several solutions.

\subsubsection{Start from the root}

The simplest solution is to use neither enhancement \sec{last changed node}
nor \sec{DF-PN-search}. In each iteration we start the search from the root.

\subsubsection{Start at the highest changed level}

We can use this solution when we have the current path in a stack, which usually
happens when we use recursion. Function \fce{updateAncestors()} will compute the level of
the highest changed node. We pop enough nodes from the stack so that the stack has
the length equal to the level of the highest changed node. In the next iteration, we can
start searching from the node on top of the stack. This way
we find the same most proving node as when we are searching from the root. We can
use the same argument as in Section \sec{last changed node} (\ref{last}) to prove it.

\subsubsection{Remember the current path}

Another variant is to start at the level of the highest changed node of the current path.
We hold the whole current path in an array. In every
node which is updated we test if it is in the current path which can be done
quickly. The level of the updated node can be found as a~number of positions
claimed by all players or we can save this information in the node.

%\section{\TODO{veta o df}}

\section{Graph representation for the clique game} 

We need to store a game state during searching the partial game tree. In other words we
need to save which edges are claimed and which color they have. In
this section, we introduce several suitable solutions and their properties.

Generally, we use an adjacency matrix but there are several ways how to
represent it in memory. An adjacency matrix is preferred over other structures
(adjacency list etc.) because most game states we search are in depth where the
colored subgraph is dense.

\subsection{Basic representation }
The adjacency matrix is straightforwardly represented as two 
dimensional array $N \times N$ of numbers where the value on the position
$(i,j)$ is zero if edge $(i,j)$ is free, or color number of the player who claimed it.

\subsection{Triangle in a~line }

The previous representation is unnecessarily large. For saving space we need just
two-bit information for every edge and we can use the symmetry of the adjacency
matrix. 

In this solution, we store the lower triangular part of the adjacency matrix in a~line of
bits. Supposing that the size of the computer word is at least ${N
\choose 2} \cdot 2$ bits we can use a single word. If this does not hold we can use more
words instead. For each edge $(i,j)$, we set two bits. First, if the edge is free
we set the bit at the position $({u \choose 2}+v)\cdot2$. Second, to store the color of
the edge we set the bit at the position $({u \choose 2}+v)\cdot2+1$.

\subsubsection{Two small improvements}

There is still space for improvement. Imagine that every two bits in a line form
a binary number -- it can be 0, 1 or 2, never 3. So, the structure from the last
paragraph is a~number in base three. We can convert this number into base two
before storing and convert it back before reading. In this case the structure has
size $ \log_2(3)/\log_2(4) \doteq 79\% $ of the original structure. It is the best solution if we care
only about space, but converting between bases is slow.

A faster possibility is to break the number into parts by three edges (stored
in six bits) and convert these short numbers into binary, which can be done
quickly using a translation table. This structure has the size $5/6 \doteq 83\%
$ of the original. A similar procedure can be used for parts of $n$ edges. If
$n=5$ the structure has size $8/10=80\%$. For a bigger $n$ the reduction in space
usage is not much better.

\subsection{Adjacency matrix in a~line}

In our algorithm, we store the whole adjacency matrix of red edges in one line of bits and blue 
in the second line of bits. Again, we suppose that we have word of size at least
$N^2$. Information if an edge has red color is stored at position $i \cdot N+j$ and
also at position $j \cdot N+i$ in the red line, and correspondingly with blue edges and
the blue line.

This structure is slightly more than twice as large than \sec{triangle in a~line}, but it has
several advantages in speed of operations, and it is easier for programming which
is also important. 

\subsubsection{A trick: Detect common neighbours}

We often need to know which nodes $v$ have red (resp. blue) edges to nodes $i$ and $j$
simultaneously. If we use basic representation there is one obvious way to find it.
We look if red edges $(i,v)$ and $(j,v)$ exist for every node $v$ ($v \ne i \;
\& \; v \ne j$). 

If we use \sec{adjacency matrix in line}, there is a better trick. We can
use binary mask and binary shift to get $i$-th and $j$-th row of the matrix, that is, bits
on positions $i \cdot N, \ldots, (i\!+\!1)\! \cdot \!N\!-\!1$ and $j \cdot N, \ldots,
(j\!+\!1)\!\cdot\!N\!-\!1 $. Then we just apply logical \fce{and} on these rows. Every nonzero bit
indicates such vertex $v$.

\subsubsection{Finding a claimed winning set} 

When the red player claims the edge $(i,j)$ we want to know if he wins (determine if he
created red $K_4$). A straightforward solution is to test every edge between every set
of four nodes which contain $i$ and $j$. 

We can also use the trick from previous subsection to do it faster. We find all
the nodes which have red edges to both $i$ and $j$. Then we test if in this
set are distinct nodes $u$ and $v$ such that red edge $(u,v)$ exists. This is illustrated
in Figure \obr{13z)}.


\subsubsection{Finding free winning sets} \label{freeK4}

It can be useful to find out if there is a~free winning set 
and how many such sets exist. We can stop searching the partial game tree if
there are no free winning sets and we can prefer to turn to edges which are
contained in most free winning sets.

Free winning sets can be found in a~similar way as the claimed winning sets. We just 
use complement of opponents line instead of players line. 

\subsubsection{Finding threats}

We want to know if the player who just played created a~threat. If he created one
threat his opponent has only one possible turn and if he created several
threat disjoint on the other two vertices, he wins. 

Threat is a~state where there are five edges with the same color and one free edge
between four nodes. Suppose that the last turn was to $(i,j)$. There are two
types of threats.
 
\begin{enumerate} 
	\item The free edge begins on node $i$ or $j$ (without loss of generality $i$). 
	\item The free edge does not contains any of them.  
\end{enumerate}

In the first case we use the trick to find all nodes $u$ which are neighbours of both
$i$ and $j$, and again to find all nodes $v$ which are neighbours of
both $j$ and $u$. If there are such $u$ and $v$, we check if there is free edge
$(i,v)$. If yes then $\{i,j,u,v\}$ is a threat. See Figure \obr{ 13a) hrozba}.

The second case is similar to finding a claimed winning set. We just check if there is a free
edge between $u,v$ instead of a red edge. See Figure \obr{ 13b) hrozba}.

\subsubsection{Computing vertex degree}

We can compute the color degree of any vertex $v$ very quickly. We precompute table of size
$2^N$ containing all the answers for degree queries. Then we just index this table by
the $v$-th row of the adjacency matrix. 

\section{Normalizing function for the clique game} \label{norm} 

Many nodes with isomorphic graphs and their subtrees are traversed during an ordinary
\sec{PN-search}. If we want to use the advantage of a high symmetry of the
game we need to join these positions and hence reduce the  number of subtrees we must
search.

In each node is a~clique game state which is an undirected complete graph with
edges colored by three colors (red, blue, free).
In an~ideal case we would search the subtree of isomorphism class
only once. When we are creating a new graph $G$ we compute a canonical
representation $H$ of the isomorphic class where $G$ belongs, and then use $H$ instead
of $G$. The enhancement \sec{transpositions to DAG} causes that the
subtree of two isomorphic graphs will be searched at most once.

However, in computing a canonical representative of a graph is a difficult
problem (see \cite{canonical}). So we compute an almost representative graph
instead. We call them a~normalise graphs.

\newtheorem*{norm}{Definition}	
\begin{norm}
{\sl A~normalising function} is a~function that for
every graph $G$ returns a~graph $H$ such that $H$ is isomorph with $G$. We we will call
$H$ a~{\sl normalise graph}.
\end{norm}

We want to find a normalise function which can be computed fast and which is
effective. Here efficiency means if we apply the function to every graph of an
isomorphism class we receive as small set of normalised graphs as possible, or
in other words there are only few distinct normalised graphs which are
isomorphic.

We introduce several such functions and compare their effectiveness and speed
on graphs\footnote{undirected 3-colored complete graphs} with five nodes and
$k$ edges colored as in clique game. See Table \ref{statsNorm}. 

\begin{table}
\centering
\begin{tabular}{c|c|c|c|c|c|c|c|c|c|c|c|c}
$k$&1&2&3&4&5&6&7&8&9&10\\
\hline
\textbf{1}&1&10&90&360&1260&2520&4200&4200&3150&1260&252\\
\textbf{2}&1&1&2&9&31&60&96&96&69&28&22\\
\textbf{3}&1&1&2&7&24&36&63&63&55&28&22\\
\textbf{4}&1&1&2&6&17&30&47&47&37&16&6\\
\end{tabular}
\caption{Efficiency of normalise functions}
\label{statsNorm}
\end{table}

\subsubsection{1 Identity}
Trivial to compute but ineffective.

\subsubsection{2 Sorting by vertex degree} \label{norm1}

We can compute the degree of each node and sort the nodes by degree. A basic improvement
is to count every red edge $M$ times where $M$ is the number of edges in a complete
graph. This improvement decreases the number of normalised graphs which are isomorphic.

\subsubsection{3 Sorting by vertex and neighbours degrees} \label{norm2}

Previous function fails in case shown in Figure \obr{14) }. In case of a tie, tt helps to
consider the degree of adjacent nodes. We define functions $f$ and $g$ on the  nodes.
\obr{14)obrazek popis: This two graphs was both normalise but isomorph before
improvement but now normalise function must return same graph.}

\begin{eqnarray*} 
f(v) = \mbox{blueDeg}(v) + \mbox{redDeg}(v) \cdot M^3.
\end{eqnarray*}

\begin{eqnarray*} 
g(v) = f(v) \cdot M^2 + \sum_{(u,v) \mbox{is red}}{f(u) \cdot M} + \sum_{(u,v) \mbox{is blue}}{f(u)}.
\end{eqnarray*}

Then the normalising function just sorts the vertices by $g$.

This function is asymptoticly slower than the previous one, but still polynomial
and more effective.

\subsubsection{4 Ideal case} 
We assign the canonical representation of graph.
This is the most effective solution but hard to compute as we noted before.
This function was not implemented so we cannot measure its speed.
Efficiency was computed in the following way: We use the previous normalising function
and for every two graphs with the same values of $g$ we straightforwardly test 
if they are isomorphic. 


\chapter{ \TODO{Kompatibilita+srovnani} enhancements }

%%%% Ukázka použití některých konstrukcí LaTeXu

\subsection{Ukázka \LaTeX{}u}
\label{ssec:ukazka}

This short subsection serves as an~example of basic \LaTeX{} constructs,
which can be useful for writing a~thesis.

Let us start with lists:

\begin{itemize}
\item The logo of Matfyz is displayed in figure~\ref{fig:mff}.
\item This is subsection~\ref{ssec:ukazka}.
\item Citing literature~\cite{lamport94}.
\end{itemize}

Different kinds of dashes:
red-black (short),
pages 16--22 (middle),
$45-44$ (minus),
and this is --- as you could have expected --- a~sentence-level dash,
which is the longest.
(Note that we have follwed \verb|a| by a~tilde instead of a~space
to avoid line breaks at that place.)

\newtheorem{theorem}{Theorem}
\newtheorem*{define}{Definition}	% Definice nečíslujeme, proto "*"

\begin{define}
A~{\sl Tree} is a connected graph with no cycles.
\end{define}

\begin{theorem}
This theorem is false.
\end{theorem}

\begin{proof}
False theorems do not have proofs.
\end{proof}

\begin{figure}
	\centering
	\includegraphics[width=30mm]{../img/logo.eps}
	\caption{Logo of MFF UK}
	\label{fig:mff}
\end{figure}


\chapter*{Conclusion}

In this thesis we introduced algorithms for solving strong positional games,
mainly PN-search and its enhancements.

We focused on games with high symmetry. The high symmetry allows, by joining
several isomorph game state into one, to narrow the partial game tree. When
partial game tree is narrow, we can search in it deeper. The other feature which
keeps the tree narrow is that we do not deal with high branching factor. For
example, we can compare our game with the game Tzaar, which has high branching
factor and low symmetry. Despite of clever heuristics and optimizations the
best bot playing the game searches the game tree only into depth
4~\cite{paulie}. For comparison, one variant of our solver for (7,4)-clique
game solved the game tree to depth 16.

There are other features distinguishing games with high symmetry. It is much
more important to join the nodes. We must treat carefully the lifetime of
the nodes, that is, we do not want to delete nodes from cache especially when their
solution was expensive because as we may find another ancestor of these nodes.
Finally in games with high symmetry much more often happens situation described
in Figure~\ref{deset}.

\subsubsection{Further work}

There are many enhancements which we did not implemented and even do
not mention. For example 1+$\epsilon$ trick.

It could also be useful to implement enhancement ``dveloping a~node without creating sons''
in Section~\ref{lastenh} and study if it can significantly reduce the partial game
tree. Also to might it may be interesting to try it on other games.

Of course, there is space for more new enhancements and heuristics. 

\addcontentsline{toc}{chapter}{Conclusion}



%%% Seznam použité literatury
%%% Seznam použité literatury je zpracován podle platných standardů. Povinnou citační
%%% normou pro bakalářskou práci je ISO 690. Jména časopisů lze uvádět zkráceně, ale jen
%%% v kodifikované podobě. Všechny použité zdroje a prameny musí být řádně citovány.

\def\bibname{Bibliography}
\begin{thebibliography}{99}
\addcontentsline{toc}{chapter}{\bibname}

%\bibitem{lamport94}
%  {\sc Lamport,} Leslie.
%  \emph{\LaTeX: A Document Preparation System}.
%  2. vydání.
%  Massachusetts: Addison Wesley, 1994.
%  ISBN 0-201-52983-1.

\bibitem{lamport94}
  {\sc Lamport,} Leslie.
  \emph{\LaTeX: A Document Preparation System}.
  2. vydání.
  Massachusetts: Addison Wesley, 1994.
  ISBN 0-201-52983-1.

\end{thebibliography}


%%% Tabulky v bakalářské práci, existují-li.
\chapwithtoc{List of Tables}

%%% Použité zkratky v bakalářské práci, existují-li, včetně jejich vysvětlení.
%\chapwithtoc{List of Abbreviations}

%%% Přílohy k bakalářské práci, existují-li (různé dodatky jako výpisy programů,
%%% diagramy apod.). Každá příloha musí být alespoň jednou odkazována z vlastního
%%% textu práce. Přílohy se číslují.
%\chapwithtoc{Attachments}

\openright
\end{document}
